\subsection{Phänomene und Eigenschaften}
\begin{sectionbox}
    \subsubsection{Grundlegende Formeln}
    \begin{description}
        \item[Ebene Welle] $\Psi(\vec r, t) = Ce^{i(k\vec r - \omega t)}$
        \item[Energie eines Quants] Kleinstmögliche Energiemenge, die ein Teilchen aufnehmen/abgeben kann $E=\hbar\omega=hf$
        \item[Impuls] $p=\hbar k$
        \item[Wellenlänge] $\lambda = \frac{2\pi}{k}=\frac{h}{mv}$. Ordnet jedem Teilchen eine Wellenlänge zu
        \item[Wellenpaket] Mehrere passende Wellen, die sich konstruktiv/destruktiv überlagern und zu einer größeren Welle interferieren
    \end{description}
\end{sectionbox}

\begin{sectionbox}
    \begin{itemize}
        \item Verhalten von Teilchen nur durch Wahrscheinlichkeiten vorhersagbar
        \item Vorhergesagte Verteilungen sind wellenartig
        \item Vorhersagen für Messgrößen nur als Mittelwert vieler Einzelmessungen bestimmbar
    \end{itemize}
\end{sectionbox}

\begin{sectionbox}
    \subsubsection{Dispersionsrelation}
    für freie Teilchen der Masse $m$ im nicht-relativistischen Fall
    \begin{equation*}
        E=\frac{p^2}{2m}=\hbar\omega \Rightarrow \omega(k) = \frac{\hbar k^2}{2m}
    \end{equation*}

    \textbf{Phasengeschwindigkeit:} Phase der ebenen Wellengleichung konstant  $\rightarrow k\vec r - \omega t = \varphi_0$,
    \begin{equation*}
        v_{\ir ph} = \left.\frac{\diff \vec r}{\diff t}\right|_{\varphi = \text{const.}}=\frac{\omega}{k}=\frac{v}{2}
    \end{equation*}

    \textbf{Gruppengeschwindigkeit eines Wellenpakets:} Beschreibt die Bewegung des Wellenpaket als ganzes (nicht die Bewegungen der einzelnen Wellen im Paket)
    \begin{equation*}
        v_{\ir gr} = \frac{\diff \omega}{\diff t}=v
    \end{equation*}
\end{sectionbox}

\begin{sectionbox}
    \subsubsection{Heisenbergsche Unschärferelation}
    Je genauer man den Ort eines Teilchens bestimmt, desto ungenauer kann man seinen Impuls bestimmen (analog andersrum)
    \begin{equation*}
        \Delta x \Delta p \ge \frac{\hbar}{2}
    \end{equation*}
\end{sectionbox}

\subsection{Wellenfunktion}
\begin{sectionbox}
    \begin{itemize}
        \item Wahrscheinlichkeitsamplitude für alle Punkte in Raum und Zeit $\Psi(\vec r, t)$
        \item Wahrscheinlichkeitsdichte gegeben durch $|\Psi(\vec r)|^2 \Rightarrow \int_V |\Psi(\vec r)|^2 \diff \vec r = 1$
    \end{itemize}
\end{sectionbox}
